\section{Appendix} \label{appendix A}

	\subsubsection*{Proofs for consequences of probability measure axioms}
	If you have not seen mathematical proofs before, then do not worry if these proofs seem complicated! It can be a 	        good idea to write down the steps yourself as you read through them.
	\subsubsection*{Proof of consequence 1}
	Let $A,B,C$ be mutually disjoint events, and suppose $P$ is a probability measure. \\
	
	 Then $A \cap C = \varnothing$ and $B \cap C = \varnothing$ (by the definition of mutually disjoint). \\
	
	\noindent So, as the union of two empty sets is also an empty set, we have $(A \cap C) \cup (B \cap C) = 	    	 
	\varnothing$. $(\star)$ \\
	
	\noindent Now, by the distributive law (see page~\pageref{distributive law}) we have $(A \cap C) \cup (B \cap C) = 
	(A \cup B) \cap C.$ \\ So, using this result along with ($\star$), we now know that $(A \cup B) \cap C = \varnothing$.  \\
	
	\noindent This means that the set $(A \cup B)$ and the set $C$ are disjoint. \\
	
	\noindent Hence by axiom 3 we have $P((A \cup B) \cup C) = P(A \cup B) + P(C).$ \;\;\; ($\dagger$) \\
	
	\noindent Also, $P(A \cup B) = P(A) + P(B)$ since $A$ and $B$ are disjoint. \\
	
	\noindent Hence, combining this last result with $(\dagger)$ we get $P(A \cup B \cup C) = P(A) + P(B) + P(C),$ \\ as 
	required. \\ 
	
	\noindent We left out the brackets around $A \cup B$ in the last step because union is associative (see 
	page~\pageref{distributive law}). \\
	
	\noindent The proof that this consequence extends to arbitrarily many events makes use of a mathematical 
	technique called ``\textit{proof by induction}'' that is beyond the scope of this course.
	
	
	\subsubsection*{Proof of consequence 2}
	We wish to prove that $P(A^c) = 1 - P(A)$. \\
	
	\noindent Firstly, recall that $\Omega = A \cup A^c$ and $A \cap A^c = \varnothing$ (see page 
	\pageref{distributive law}). \\
	
	\noindent By axiom 1 we have $P(\Omega) = 1$. Thus, as $A \cup A^c = \Omega$, we have $P(A \cup A^c) = 1$ 
	also.  $(\star)$ \\
	
	\noindent Also, since $A \cap A^c = \varnothing$ means that $A$ and $A^c$ are disjoint, we have by axiom 3 that 
	\\ $P(A \cup A^c) = P(A) + P(A^c)$. \\
	
	\noindent Combining this last result with ($\star$) we get $P(A) + P(A^c) = 1$. \\ Rearranging gives $P(A^c) = 1 - 
	P(A)$, as required.
	
	\subsubsection*{Proof of consequence 3}
	We wish to prove that $P(\varnothing) = 0$. This follows easily from consequence 2, and the fact that $\varnothing 
	= \Omega^c$. \\
	
	We have, using consequence 2, $P(\varnothing) = P( \Omega^c) = 1 - P(\Omega).$ Now by axiom 1 we have 
	$P(\Omega) = 1$ so $P(\varnothing) = 1 - 1 = 0,$ as required.
	
	\subsubsection*{Proof of consequence 4}
	Let $A,B$ be two events such that $A \subseteq B.$ \\
	
	\noindent We wish to prove that $P(A) \leq P(B)$. \\
	
	
	\noindent Notice that since $A$ is fully contained within $B$, we must have $B = A \cup C$ for some 
	(possibly~empty) event~$C$. We can choose $C$ so that $A \cap C = \varnothing$. \\
	
	\noindent Then by axiom 3, $P(B) = P(A \cup C) = P(A) + P(C)$ (since we chose $C$ so that $A$ and $C$ are 
	disjoint). \\
	
	\noindent Now, by axiom 2 we have~$P(C)~\geq~0$. So $P(B)$ is equal to $P(A)$ plus some number greater than 
	or equal to zero. Hence $P(B) \geq P(A),$ as required.
	
	\subsubsection*{Proof of consequence 5}
	To get an intuitive idea of why this consequence is true, consider the fact from page \pageref{distributive law}, 
	that Size of $A\cup B$ = Size of $A$ + Size of $B$ - Size of $A \cap B$. The actual proof requires some ideas about 
	sets that have not yet been introduced in this course, and it is slightly outside of the focus of the course. Hence 
	some readers may wish to skip this proof. \\
	
	\noindent Let A,B be (not necessarily disjoint) sets, and suppose $P$ is a probability measure. We define the set 
	$A\setminus B$ to be the set obtained from $A$ by excluding all elements that also belong to $B$ (i.e. by 
	excluding all elements from $A \cap B$). We define $B\setminus A$ analogously. \\
	
	\noindent Notice this implies that $A =(A\setminus B) \cup (A \cap B)$ and similarly $B = (B\setminus A) \cup (A 
	\cap B)$. Furthermore it implies that $(A\setminus B) \cap (A \cap B) = \varnothing$ and $(B\setminus A) \cap (A 
	\cap B) = \varnothing$. \\
	
	\noindent Hence by axiom 3 we have $P(A) = P(A\setminus B) + P(A \cap B)$ and $P(B) = P(B\setminus A) + P(A 
	\cap B)$. \\
	
	\noindent So $P(A) + P(B) = P(A\setminus B) + P(B\setminus A) + 2P(A \cap B)$. \qquad $(\star)$ \\
	
	\noindent With some consideration we can also see that $(A\setminus B) \cap (B\setminus A) = \varnothing$. \\ 
	So the sets~$(A\setminus B),~(B\setminus A)$,~and~$A \cap B$ are mutually disjoint. Hence by consequence 1 
	we have $P(A\setminus B) + P(B\setminus A) + P(A \cap B) = P((A\setminus B) \cup (B\setminus A) \cup (A \cap B)) 
	= P(A \cup B)$. \\
	
	\noindent Therefore, substituting this result into $(\star)$ gives \\
	$P(A) + P(B) = P(A \cup B) + P(A \cap B)$. \\
	
	So  $P(A \cup B) = P(A) + P(B) - P(A \cap B)$, as required.
